\section{Introduction}
\label{sec:intro}

Autoregressive and streaming video diffusion models now extend generation from seconds to minute-scale horizons. Evaluation, however, has not kept pace: small temporal errors that are invisible in short clips accumulate over long rollouts into background displacement, color drift, repeated textures, or motion decay, and whole-frame metrics summarize these qualitatively different failures into a single score.

The problem is sharpest in \textbf{fixed-camera nature-flow generation}. In a static-camera video of a river, the water should move while the riverbank, rocks, and surrounding landscape remain stationary. The same separation governs waterfalls, ocean waves, rainfall, fire, smoke, and wind-driven vegetation: some regions form \emph{static support} that must stay structurally fixed, while others carry \emph{dynamic flow} that must persist. Whole-frame motion measures reward optical flow regardless of its source---crediting background drift as readily as genuine fluid motion---while temporal-consistency measures can favor videos whose stability comes from suppressing the very flow they should sustain. Two opposite failures, \textbf{static drift} and \textbf{flow decay}, are thereby conflated, and increasingly so at long horizons, where accumulated background motion can dominate a sequence's apparent dynamics.

This ambiguity spans conditioning modalities. In text-to-video (T2V) generation, the model must establish static support and preserve it without a visual reference; in image-to-video (I2V) generation, the input image defines that support explicitly. Both tasks pose the same fixed-camera question: \emph{does motion persist where it is intended, without leaking into regions that should remain static?}

We introduce \textbf{SNF-Bench} (Static--Natural Flow), an evaluation framework built around this static--dynamic decomposition, with category-balanced T2V and I2V tracks evaluated primarily at 60- and 120-second horizons (shorter and ultra-long horizons serve as diagnostics). Evaluation is organized into three factors. \textbf{Static Fidelity} measures geometric displacement and residual optical flow inside nominally stationary regions; for I2V, it additionally compares static regions against the shared input image. \textbf{Flow Persistence} jointly reports the retention of dynamic-region motion over time and its absolute drift-compensated magnitude, so a persistence ratio cannot look strong merely because motion is uniformly weak. \textbf{Drift Leakage} estimates how much apparent dynamic-region motion is explained by global background displacement.

A benchmark is only as trustworthy as its measurements, so we validate the protocol independently of any model comparison. Controlled perturbations inject known translation, rotation, scale and color drift, motion attenuation, temporal repetition, and generation-like corruptions into fixed-camera videos, testing that each metric responds selectively and monotonically to the failure it targets. We calibrate the drift-leakage estimate on sequences with known mixtures of global displacement and local motion, quantify region-mask agreement, and verify that conclusions are stable under mask erosion and dilation.

We then audit public long-horizon generators under two explicitly separated settings: \textbf{native-capability evaluation}, which runs each method in its intended configuration and records its supported horizon, block size, denoising schedule, and memory policy; and a secondary \textbf{matched stress test} under a common fixed-camera protocol. The analysis reveals cases where conventional metrics and SNF-Bench interpret the same outputs differently---high whole-frame motion coexisting with substantial estimated drift leakage, and high stability coexisting with decaying flow. We do not argue that existing metrics are wrong; we argue that whole-frame aggregation does not represent a spatial distinction that this setting requires. Nor does SNF-Bench measure physical realism: a video may be drift-free and persistent yet physically implausible. Our goal is narrower and complementary---making the spatial origin and temporal persistence of motion explicit.

Our contributions are:
\begin{enumerate}
\item \textbf{SNF-Bench}, a controlled evaluation framework for long-horizon fixed-camera nature-flow generation, with T2V and I2V tracks unified by a shared decomposition of static support and dynamic flow.
\item A \textbf{drift--motion evaluation protocol} separately measuring static fidelity, flow persistence with absolute magnitude, and drift leakage.
\item \textbf{Controlled validation} of every measurement via synthetic and generation-like perturbations, including drift-leakage calibration and mask-robustness analysis.
\item A \textbf{public-model audit} under native and matched configurations, showing that the fixed-camera decomposition materially changes the interpretation and ranking of long-horizon behavior relative to generic metrics.
\end{enumerate}